\documentclass{article}
\usepackage{bm}
\usepackage{mathtools}
\usepackage{amsthm}
\usepackage{fancyhdr}
\usepackage{float}
\usepackage{listings}
\lstset{breaklines=true}
\usepackage{graphicx}
\usepackage{xcolor}
\usepackage{titlesec}
\usepackage{titling}
\usepackage{tabularx}
\usepackage{enumitem}
\usepackage{amsmath}
\usepackage{amsfonts}
\usepackage{hyperref}
\usepackage[margin=1in]{geometry}

\hypersetup{
	colorlinks=true,
	linkcolor=blue,
	filecolor=magenta,
	urlcolor=blue,
	citecolor=blue,
	anchorcolor=blue
}

\renewcommand{\thesubsubsection}{\Roman{subsubsection}}
\title{Math 3IA3 - Assignment 3}
\author{Matthew Farah, \href{mailto:farahm11@mcmaster.ca}{$\langle$farahm11@mcmaster.ca$\rangle$}, 400468251}
\date{\today}

\makeatletter
\setlist{listparindent=\parindent}

\pagestyle{fancy}
\fancyhead{}
\fancyhead[L]{Matthew Farah, $\langle$\href{mailto:farahm11@mcmaster.ca}{farahm11@mcmaster.ca}$\rangle$, 400468251}
\fancyhead[R]{Math 3IA3 - Assignment 3}

\newcolumntype{Y}{>{\centering\arraybackslash}X}
\setlength{\parindent}{0cm}

\setcounter{section}{3}

\newcounter{chapter}
\setcounter{chapter}{4}

\newcounter{exercise}
\renewcommand{\theexercise}{\arabic{exercise}}

\newenvironment{exercise}[1][]{
	\ifx\\#1\\
		\refstepcounter{exercise}
	\else
		\setcounter{exercise}{#1}
	\fi
	\addcontentsline{toc}{section}{\protect{\large{Exercise \thesection.\thechapter.\theexercise}}}
	\large{\par\noindent\textbf{Exercise \thesection.\thechapter.\theexercise}}\vspace{.5em}\par\noindent
}{\vspace{1.5em}}

\newenvironment{solution}{
	\vspace{0.4cm}\par\noindent\large\textbf{{Solution:}}\par\vspace{1em}
}{
	\vspace{2.5em}
	\newpage
}

\newcounter{partslist}

\makeatletter
\newcommand{\parts@seriesname}{parts@\theexercise}

\newenvironment{parts}{
	\edef\parts@ser{\parts@seriesname}
	\ifcsname\parts@ser\endcsname
		\begin{enumerate}[resume=\parts@ser,label=\textbf{(\alph*)},leftmargin=2.0em,itemsep=0.4em,before=\large]
	\else
		\begin{enumerate}[series=\parts@ser,label=\textbf{(\alph*)},leftmargin=2.0em,itemsep=0.4em,before=\large]
		\expandafter\gdef\csname\parts@ser\endcsname{1}
	\fi
}{
	\end{enumerate}
}

\makeatother

\makeatletter

\newcounter{currentstep}[partslist]
\renewcommand{\thecurrentstep}{\Roman{currentstep}}

\newenvironment{steps}{
	\edef\steps@ser{steps@\theexercise}
	\ifcsname\steps@ser\endcsname
		\begin{enumerate}[
			resume=\steps@ser,
			label=\bfseries{(\Roman*)},
			leftmargin=1.5em,
			itemsep=0.4em,
			topsep=0.8em
		]
	\else
		\begin{enumerate}[
			series=\steps@ser,
			label=\bfseries{(\Roman*)},
			leftmargin=1.5em,
			itemsep=0.4em,
			topsep=0.8em
		]
		\expandafter\gdef\csname\steps@ser\endcsname{1}
	\fi
	\let\steps@saved@item\item
	\renewcommand{\item}{\stepcounter{currentstep}\steps@saved@item}
}{
	\let\item\steps@saved@item
	\end{enumerate}
}

\makeatother

\begin{document}

\maketitle
\pagestyle{fancy}

\tableofcontents
\newpage

\begin{exercise}[6]
Let $x_n \coloneq n^{\frac{1}{n}}$ for $n \in \mathbb{N}$.
\end{exercise}

\begin{parts}
	\item Show that $x_{n + 1} < x_n$ if and only if ${(1 + \frac{1}{n})}^{n} < n$, and infer that the inequality is valid for $n \ge 3$. Conclude that $(x_n)$ is ultimately decreasing and that $x \coloneq \lim_{n\to\infty}(x_n)$ exists.
\end{parts}

\begin{solution}

	\begin{steps}
		\item Show that $x_{n + 1} < x_n \iff (1 + \frac{1}{n})^n < n$.
	\end{steps}

	Using basic algebraic manipulations, we find:

	\begin{align*}
		x_{n + 1} &< x_n \\
		\Updownarrow \\
		(n + 1)^{\frac{1}{n + 1}} &< n^{\frac{1}{n}} \\
		\Updownarrow \\
		((n + 1)^{\frac{1}{n + 1}})^{n + 1} &< (n^{\frac{1}{n}})^{n + 1} \\
		\Updownarrow \\
		n + 1 &< (n^{\frac{1}{n}})^{n} \cdot n^{\frac{1}{n}} \\
		\Updownarrow \\
		n + 1 &< n \cdot n^{\frac{1}{n}} \\
		\Updownarrow \\
		\frac{n + 1}{n} &< n^{\frac{1}{n}} \\
		\Updownarrow \\
		1 + \frac{1}{n} &< n^{\frac{1}{n}} \\
		\Updownarrow \\
		\left( 1 + \frac{1}{n} \right)^n &< (n^{\frac{1}{n}})^n \\
		\Updownarrow \\
		\left( 1 + \frac{1}{n} \right)^n &< n \\
	\end{align*}

	Therefore, $x_{n + 1} < x_n \iff (1 + \frac{1}{n})^n < n$

	\begin{steps}
		\item Show that $(1 + \frac{1}{n})^n < n, \; \forall n \ge 3$.
	\end{steps}

	We will proceed to prove $(1 + \frac{1}{n})^n < n, \; \forall n \ge 3$ via induction. Let $P(n)$ be the statement $P(n): (1 + \frac{1}{n})^n < n$.

	\begin{center}
		\underline{Base Case:}
	\end{center}

	\begin{align*}
		&\text{LHS:} &\text{RHS} \\
		&=(1 + \frac{1}{n})^n &=n \\
		&= \left( 1 + \frac{1}{3} \right)^3 &=3 \\
		&= \left( \frac{4}{3} \right)^3 &=3 \\
		&=\frac{64}{27} &=3 \\
	\end{align*}

	Therefore, since LHS $<$ RHS, the base case holds and $P(3)$ is true.

	\begin{center}
		\underline{State Induction Hypothesis:}
	\end{center}

	Let $k \in \mathbb{N}$ be arbitrary, we can therefore formulate our induction hypothesis as:

	\begin{align*}
		P(k): (1 + \frac{1}{k})^k < k \\
	\end{align*}

	\begin{center}
		\underline{Prove $P(k) \implies P(k + 1)$:}
	\end{center}

	\begin{align*}
		(1 + \frac{1}{k + 1})^{k + 1} &= (1 + \frac{1}{k + 1})^{k} \cdot (1 + \frac{1}{k + 1}) \\
		&\le (1 + \frac{1}{k})^{k} \cdot (1 + \frac{1}{k}) \qquad \langle \; \text{Since} \; \frac{1}{k + 1} < \frac{1}{k} \; \forall k \in \mathbb{N} \; \rangle \\
		&< k \cdot (1 + \frac{1}{1 + k}) \qquad \langle \; \text{By induction hypothesis} \; \rangle \\
		&= k + 1
	\end{align*}

	Therefore, $P(k) \implies P(k + 1)$.

	Since $P(k) \implies P(k + 1)$ for arbitrary $k \in \mathbb{N}$, and since we know $P(3)$ is true, then by principle of mathematical induction, we know that $P(n)$ holds $\forall n \in \mathbb{N}$ such that $n \ge 3$.

	\begin{steps}
		\item Show $x \coloneq \lim_{n\to\infty}(x_n)$ exists.
	\end{steps}

	From parts (I) and (II), since $x_{n + 1} < x_{n} \iff \left( 1 + \frac{1}{n} \right)^n < n$, and $\left( 1 + \frac{1}{n} \right)^n < n$ holds for all $n \ge 3$, we have that $x_{n + 1} < x_{n}$ is true $\forall n \ge 3$ and thus the sequence $x_n$ is eventually monotone. Furthermore, since $n \ge 1$, we know that $(n)^\frac{1}{n} \ge 1$ and thus, since $(x_n)$ is monotonously decreasing past $n = 3$ and is lowerbounded by 0, by the Monotone Convergence Theorem, we know that $(x_n)$ converges.
\end{solution}

\begin{parts}
	\item Use the fact that the subsequence $(x_{2n})$ also converges to x to conclude that $x = 1$.
\end{parts}

\begin{solution}
	Since we know that $(x_n)$ converges, we know that all subsequences of $(x_n)$ also converge to $x$. Therefore, using the limit identity proven in the textbook that $\lim_{n\to\infty}(p)^\frac{1}{n} = 1 \; \forall p \in \mathbb{R}, r \ge 1$, we can find that:
	
	\begin{align*}
		\lim_{n\to\infty} \left( (2n)^{\frac{1}{2n}} \right) &= x \implies \\
		\lim_{n\to\infty} \left( (2n)^{\frac{1}{n}} \right) &= x^2 \implies \\
		\lim_{n\to\infty} \left( 2^{\frac{1}{n}} n^{\frac{1}{n}} \right) &= x^2 \implies \\
		\lim_{n\to\infty} \left( 2^{\frac{1}{n}} \right) \cdot \lim_{n\to\infty} \left( n^{\frac{1}{n}} \right) &= x^2 \implies \\
		1 \cdot x &= x^2 \implies \\
		x &= x^2 \implies \\
		x^2 - x &= 0 \implies \\
		x(x - 1) &= 0 \\
	\end{align*}

	Therefore, $x = 1$ or $x = 0$, but since we know that $x_n \ge 1 \; \forall n \in \mathbb{N}$, $x$ must equal 1. Therefore, $\lim_{n\to\infty}(x_n) = 1$.
\end{solution}

\setcounter{chapter}{5}

\begin{exercise}[5]
	If $x_n \coloneq \sqrt{n}$, show that $(x_n)$ satisfies $\lim_{n\to\infty}|x_{n + 1} - x_n| = 0$ but is not a Cauchy sequence.
\end{exercise}

\begin{solution}

	\begin{steps}
		\item Show that $x_n$ satisfies $\lim_{n\to\infty}|x_{n + 1} - x_n| = 0$.
	\end{steps}
	
	Taking it as given that the sequence $\frac{1}{n}$ converges to zero, by algebraic limit rules, we can determine that the sequence $\frac{1}{\sqrt{n}}$ also converges to zero. Thus, for arbitrary $\epsilon > 0$, we know that there exists some $N \in \mathbb{N}$ such that $\forall \ge N, \frac{1}{\sqrt{n}} < \epsilon$. Therefore, for all such $n$, we have that:

	\begin{align*}
		|(x_{n + 1} - x_n) - 0| &= |\sqrt{n + 1} - \sqrt{n}| \\\\
		&= \sqrt{n + 1} - \sqrt{n} \qquad \langle \; \text{Since} \sqrt{n + 1} > \sqrt{n} \; \forall n \in \mathbb{N} \; \rangle \\\\
		&= \sqrt{n + 1} - \sqrt{n} \cdot \frac{\sqrt{n + 1} + \sqrt{n}}{\sqrt{n + 1} + \sqrt{n}} \\
		&= \frac{\left( \sqrt{n + 1} - \sqrt{n} \right) \left ( \sqrt{n + 1} + \sqrt{n} \right)}{\sqrt{n + 1} + \sqrt{n}} \\
		&= \frac{n + 1 - n}{\sqrt{n + 1} + \sqrt{n}} \\
		&= \frac{1}{\sqrt{n + 1} + \sqrt{n}} \\
		&\le \frac{1}{\sqrt{n}} \\
		&< \epsilon \\
	\end{align*}

	Thus, since for arbitrary $\epsilon > 0$, we can find an $N \in \mathbb{N}$ such that for all $n \ge N$, $|x_{n + 1} - x_n| < \epsilon$, by definition of convergence, $\lim_{n\to\infty}|x_{n + 1} - x_n| = 0$.

	\begin{steps}
		\item Show that $x_n$ is not a Cauchy sequence.
	\end{steps}

	We show that $x_n$ is not a Cauchy sequence by proving that it diverges because, over the set of real numbers, a sequence is Cauchy if and only if it is convergent to a real number. Let $M^2 > 0$ be arbitrary. By the Archimedean property, we know that there exists an $N \in \mathbb{N}$ such that for all $n \ge N \implies n > M^2$. Thus:

	\begin{align*}
		n &\ge N \implies \\
		n &> M^2 \implies \\
		\sqrt{n} &> \sqrt{M^2} \implies \\
		x_n &> M \\
	\end{align*}

	Therefore, we know that $\lim_{n\to\infty} (x_n)$ diverges to infinity and therefore $(x_n)$ cannot be Cauchy.
\end{solution}

\setcounter{chapter}{6}

\begin{exercise}[10]
	Show that if $\lim_{n\to\infty}(\frac{a_n}{n}) = L$, where $L > 0$, then $\lim_{n\to\infty}(a_n) = +\infty$.
\end{exercise}

\begin{solution}
	Let $M > 0$ be arbitrary. Since $\lim_{n\to\infty} \frac{1}{n} = 0$ is known, we can conclude using algebraic limit rules that $\lim_{n\to\infty}(\frac{M}{N}) = 0$. Since $\lim_{n\to\infty} (\frac{a_n}{n}) = L > 0$, we know there exists some $N_0 \in \mathbb{N}$ such that $\forall n \ge N_0$:

	\begin{gather*}
		|\frac{a_{n}}{n} - L| < \frac{L}{2} \\
	\end{gather*}

	And therefore, $\frac{L}{2} < \frac{a_n}{n} < \frac{3}{2}L$ for all such $n$. Likewise, we know that there exists some $N_1 \in \mathbb{N}$ such that $\forall n \ge N_1$:

	\begin{align*}
		|\frac{M}{n}| &< \frac{L}{2}
	\end{align*}

	And therefore, $0 < \frac{M}{n} < \frac{L}{2}$ for all such n. Thus, letting $N = \text{max}\{N_0, N_1\}$, we have that $\forall n \ge N$:

	\begin{align*}
		\frac{a_n}{n} &> \frac{M}{n} \\
		a_n &> M
	\end{align*}

	Therefore, since for arbitrary $M > 0$, we can always find an $N \in \mathbb{N}$ such that $a_n > M \; \forall n \ge N$, $\lim_{n\to\infty}(a_n) = +\infty$.
\end{solution}

\setcounter{chapter}{7}

\begin{exercise}[11]
	If $\sum_{n = 1}^{\infty}a_n$ with $a_n > 0$ is convergent, then is $\sum_{n = 1}^{\infty}(a_n)^2$ always convergent? Either prove it or give a counter-example.
\end{exercise} 

\begin{solution}
	\begin{proof}
		We know, by definition of infinite sums, that if $\sum_{n = 1}^{\infty}a_n$ convergent to some limit $a$, then the sequence of its partial sums, defined by $(s_k) = \sum_{n = 1}^{k}(a_n)$ must converge likewise converge to $a$. Therefore, we know that $\lim_{k\to\infty}(s_k) = a$. Using the same logic, we know that the infinite sum $\sum_{n = 1}^{\infty}(a_n)^2$ converges if the sequence of its partial sums converges.

		Since $a_n > 0, \; \forall n \in \mathbb{N}$, we know that $(s_k)$ is increasing and since $a_n > 0 \implies (a_n)^2 > 0$, we likewise know the sequence of partial sums of $(a_n)^2$ is likewise increasing. Furthermre, as a corollary to the fact $(s_k)$ is increasing, we know that $a$ must upperbound $(s_k)$. Now, algebraic limit rules dictate that $\left( \lim_{k\to\infty}(s_k) \right)^2 = a^2$, however, notice that:

		\begin{align*}
			(s_k^2) &= \left( \sum_{n = 1}^{k}(a_n) \right)^2 \\
			&= \sum_{n = 1}^{k}(a_n) \cdot \sum_{n = 1}^k(a_n) \\
			&= (a_1 + a_2 + \ldots + a_k) \cdot (a_1 + a_2 + \ldots + a_k) \\
			&= a_1^2 + a_1a_2 + \ldots + a_1a_k + a_2a_1 + a_2^2 + \ldots a_2a_k + \ldots + a_ka_1 + a_ka_2 + \ldots + a_k^2 \\
			&\ge a_1^2 + a_2^2 + \ldots + a_k^2 \qquad \langle \; \text{Since} \; a_n > 0 \; \forall n \in \mathbb{N} \; \rangle \\
			&= \sum_{n = 1}^k(a_n)^2 \\
		\end{align*}

		Therefore, we know that the sequence of partial sums of $0 < \sum_{n = 1}^k(a_n)^2 \le s_k$, we therefore know that $a$ must also upperbound the sequence of partial sums of $(a_k)^2$. Thus, since the sequence of partial sums of $(a_k)^2$ is both increasing and upperbounded by $a$, by Monotone Convergence Theorem, the sequence converges and thus, the infinite sum $\sum_{n = 1}^{\infty}(a_n)^2$ converges.
	\end{proof}
\end{solution}

\setcounter{section}{4}
\setcounter{chapter}{1}

\begin{exercise}[7]
	Show that $\lim_{x\to c}(x^3) = c^3$ for any $c \in \mathbb{R}$.
\end{exercise}

\begin{solution}
	We can easily show this using the sequential definition of functional limits. Let $(x_n)$ by any sequence defined on $\mathbb{R}$ (since the domain of $x^3$ is $\mathbb{R}$) such that $x_n \neq c \; \forall n \in \mathbb{N}$ and $\lim_{n\to\infty}(x_n) = c$. By the algebraic limit theorems, we know that:

	\begin{align*}
		\lim_{n\to\infty}(x^3) &= (\lim_{n\to\infty}(x_n))^3 \\
		&= c^3 \\
	\end{align*}

	Therefore, by definition of functional limits, since for all such $(x_n)$, $\lim_{n\to\infty}(x_n)^3 = c^3$, $\lim_{x \to c}(x)^3 = c^3$.
\end{solution}
\end{document}
